\documentclass[11pt]{article}

\usepackage[margin=1in]{geometry}
\usepackage{amsmath, amssymb}
\usepackage{xcolor}
\usepackage{listings}
\usepackage{hyperref}
\usepackage{array}
\usepackage{booktabs}

\hypersetup{
    colorlinks=true,
    linkcolor=blue!60!black,
    urlcolor=blue!60!black,
    citecolor=blue!60!black
}

\lstdefinestyle{bugpython}{
  language=Python,
  basicstyle=\ttfamily\small,
  keywordstyle=\color{blue!60!black},
  stringstyle=\color{green!40!black},
  commentstyle=\color{gray!70!black},
  showstringspaces=false,
  breaklines=true,
  frame=single,
  rulecolor=\color{gray!30},
  columns=fullflexible,
  keepspaces=true
}

\title{SymPy Bug Report}
\date{}

\begin{document}

\author{}
\maketitle

\section{Bug 37: Wrong Sign in an Antiderivative on the Negative Real Branch}

\subsection{Minimal Reproducer}

\medskip

\begin{lstlisting}[style=bugpython]
from sympy import N, diff, integrate, sqrt, symbols

x = symbols("x")
integrand = 1/(x*sqrt(x**2 - 1))
F = integrate(integrand, x)
dF = diff(F, x)

print("antiderivative =", F)
print("derivative =", dF)
print("integrand at x=-3:", N(integrand.subs(x, -3), 50))
print("derivative at x=-3:", N(dF.subs(x, -3), 50))
print("difference:", N((dF - integrand).subs(x, -3), 50))
\end{lstlisting}

\subsection{Observed Behavior}

\medskip

\begin{lstlisting}[style=bugpython]
antiderivative = Piecewise((I*acosh(1/x), 1/Abs(x**2) > 1), (-asin(1/x), True))
derivative = Piecewise((-I/(x**2*sqrt(-1 + 1/x)*sqrt(1 + 1/x)), 1/Abs(x**2) > 1), (1/(x**2*sqrt(1 - 1/x**2)), True))
integrand at x=-3: -0.11785113019775792073347406035080817321413932294808
derivative at x=-3: 0.11785113019775792073347406035080817321413932294808
difference: 0.23570226039551584146694812070161634642827864589616
\end{lstlisting}

The returned expression differentiates to the wrong real sign at the ordinary
point \(x=-3\).  The residual is not zero: the derivative of SymPy's returned
antiderivative is the positive opposite of the requested integrand on this
negative real branch.

\subsection{Expected Behavior}

For an indefinite integral returned as an antiderivative, differentiating the
result should recover
\[
    \dfrac{1}{x\sqrt{x^2-1}}
\]
at ordinary points of the selected real branch.  In particular, at \(x=-3\) the
derivative should evaluate to
\[
    -0.11785113019775792073347406035080817321413932294808\ldots,
\]
matching the integrand rather than its negation.

\subsection{Mathematical Explanation}

On the real interval \(x<-1\), the principal square root satisfies
\(\sqrt{x^2-1}>0\), while \(x<0\).  Therefore
\[
    \dfrac{1}{x\sqrt{x^2-1}} < 0 .
\]
SymPy's fallback branch is \(-\operatorname{asin}(1/x)\).  Differentiating it gives
\[
    \dfrac{d}{dx}\left[-\operatorname{asin}\left(\dfrac{1}{x}\right)\right]
    =
    \dfrac{1}{x^2\sqrt{1-1/x^2}} .
\]
For \(x<-1\), this derivative is positive, because both \(x^2\) and
\(\sqrt{1-1/x^2}\) are positive.  It is related to the target integrand by
\[
    \dfrac{1}{x^2\sqrt{1-1/x^2}}
    =
    \dfrac{1}{|x|\sqrt{x^2-1}},
\]
which equals the requested integrand on \(x>1\) but equals its negation on
\(x<-1\).  The discrepancy is therefore not a harmless additive constant,
simplification choice, or equivalent branch representation: the derivative has
the wrong sign on a nonsingular real interval.

\subsection{Source-Level Diagnosis}

The diagnosis narrows the source of the bad result to the Meijer-G indefinite
integration path in \nolinkurl{sympy/integrals/meijerint.py}, around lines
1653--1777.  The relevant functions named in the diagnosis are:

\begin{lstlisting}[style=bugpython]
meijerint_indefinite
_meijerint_indefinite_1
\end{lstlisting}

The public call to \texttt{integrate} constructs an \texttt{Integral} and reaches
\nolinkurl{Integral._eval_integral} in
\nolinkurl{sympy/integrals/integrals.py}.  The ordinary Risch and heuristic
paths do not handle this algebraic integrand, so the Meijer-G fallback calls
\texttt{meijerint\_indefinite}; direct probing of that function reportedly
produces the same bad \texttt{Piecewise} result.

The diagnosed path rewrites the integrand as a Meijer-G expression, applies a
generic antiderivative formula, expands the result with \texttt{hyperexpand},
combines powers using \texttt{powdenest(..., polar=True)}, and then removes
polar or branch annotations with \texttt{\_my\_unpolarify(\_clean(...))}.  The
important outcome is that the final fallback branch is \texttt{-asin(1/x)} under
a \texttt{True} condition, after an inner branch guarded by
\texttt{1/Abs(x**2) > 1}.  That fallback conflates the two exterior real
intervals \(x>1\) and \(x<-1\).

The source-level behavior causes the observed mathematical failure because
\(-\operatorname{asin}(1/x)\) is a correct real antiderivative on the positive exterior branch
but has the opposite derivative sign on the negative exterior branch.  The
diagnosis therefore identifies a branch/sign loss during the
Meijer/hyperexpanded expression's conversion back to ordinary inverse
trigonometric functions.  A plausible fix direction supported by the diagnosis
is to retain a sign-dependent \texttt{Piecewise} result for \(x>1\) and \(x<-1\),
or to decline this simplified inverse-trigonometric antiderivative when the
branch cannot be made valid over the whole fallback condition.

\subsection{Additional Failing Instances}

The same failure appears across representative negative real sample points in
the interval \(x<-1\).  The parametrized verification checks
\(x=-2,-3,-4,-5,-6,-7\).  In each case, the returned fallback branch
differentiates to \(1/(|x|\sqrt{x^2-1})\), while the requested integrand is
\(1/(x\sqrt{x^2-1})\); since \(x<0\), the two values have opposite signs.

\begin{center}
\small
\renewcommand{\arraystretch}{1.3}
\begin{tabular}{@{}>{\raggedright\arraybackslash}p{0.44\linewidth}>{\raggedright\arraybackslash}p{0.23\linewidth}>{\raggedright\arraybackslash}p{0.23\linewidth}@{}}
\toprule
\textbf{Input / Expression} & \textbf{SymPy behavior} & \textbf{Expected behavior} \\
\midrule
\(x=-2,\; F'(-2)\) & \(+1/(2\sqrt{3})\) & \(-1/(2\sqrt{3})\) \\
\(x=-3,\; F'(-3)\) & \(+1/(6\sqrt{2})\) & \(-1/(6\sqrt{2})\) \\
\(x=-4,\; F'(-4)\) & \(+1/(4\sqrt{15})\) & \(-1/(4\sqrt{15})\) \\
\(x=-5,\; F'(-5)\) & \(+1/(10\sqrt{6})\) & \(-1/(10\sqrt{6})\) \\
\(x=-6,\; F'(-6)\) & \(+1/(6\sqrt{35})\) & \(-1/(6\sqrt{35})\) \\
\(x=-7,\; F'(-7)\) & \(+1/(28\sqrt{3})\) & \(-1/(28\sqrt{3})\) \\
\bottomrule
\end{tabular}
\end{center}

\subsection{Related Issues Assessment}

The bug-finding harness performed a best-effort deduplication check and classified
this as a new instance of a known failure mode.  The closest public precedents in
the same \(\sqrt{x^2-1}\) integration family are two open/closed SymPy issues:
\href{https://github.com/sympy/sympy/issues/20982}{sympy/sympy\#20982} (closed),
``integrate(1 / (x ** 2 * sqrt(x ** 2 - 1))) gives wrong integral (lost sign when
\(x<0\))'', which reports the same negative-real sign-loss symptom on a sibling
\(1/(x^2\sqrt{x^2-1})\) integrand; and
\href{https://github.com/sympy/sympy/issues/23566}{sympy/sympy\#23566} (open),
``integrate(1/sqrt(x**2-1)) shouldn't be acosh(x)'', which concerns the same
Meijer-G \(\operatorname{acosh}/\operatorname{asin}\) \texttt{Piecewise} branch
handling for \(\sqrt{x^2-1}\).  No public issue or pull request, however, targets
the exact \(1/(x\sqrt{x^2-1})\) indefinite integral having the wrong sign on
negative real inputs.  The broad branch-cut family is therefore known and
publicly documented, while this minimized negative-branch counterexample is
specific and previously unreported, remaining individually actionable as an issue
and regression test.

\subsection{Proposed SymPy Regression Test}

\medskip

\begin{lstlisting}[style=bugpython]
import pytest
from sympy import N, diff, integrate, sqrt, symbols


@pytest.mark.parametrize("x_value", [-2, -3, -4, -5, -6, -7])
def test_integral_of_inverse_x_sqrt_x2_minus_1_negative_branch(x_value):
    x = symbols("x")
    integrand = 1/(x*sqrt(x**2 - 1))
    antiderivative = integrate(integrand, x)
    residual = diff(antiderivative, x) - integrand
    assert abs(complex(N(residual.subs(x, x_value), 50))) < 1e-40
\end{lstlisting}

\end{document}
